\section{Research Plan}

\subsection{Thrusts}

There are two major thrusts of this proposal, both including
conceptual and engineering aspects, with feedback between concept,
evaluation, and engineering guiding the research activity.

The first thrust, and the most important, is the exploration of the
ability of relational coverage concepts to improve the effectiveness
of differential fuzzing, given a fixed, well-defined, problem of
detecting behavior divergences among two or more implementations of
the same functionality.  We believe the core idea of relational
coverage is likely to significantly enhance the effectiveness of
differential fuzzing for bug detection.  This aspect of the project
will consist of implementing the simplest plausible versions of the
kinds of ideas proposed above, and exploring their effectiveness via
preliminary evaluations of the techniques using the same evaluation
techniques as proposed below for the publication of the research.  We
expect some feedback from experiments to be necessary, since adding
new core techniques to AFL++ will involved many detailed decisions
that may impact fuzzing effectiveness but will either be
underdetermined by the basic concepts or will involve fuzzer
throughput driven tradeoffs (approximations similar to those in the
standard fuzzer compression of coverage information, for example).

The second thrust involves two sub-parts: first, since many potential users of
differential fuzzing are dissuaded from applying the technique by the
cost of constructing a differential harness, we aim to develop a
principled approach to automatic calibration of harnesses using a
variety of possible difference signals (output to stdout/stderr, file
outputs, time and space performance, etc.).  Second, because some of
the most interesting and important differential fuzzing problems
require careful thought and complex abstraction mechanisms, or
removing parts of the input space not handled by one or more
differentially compared implementations, we will investigate the
development of a DSL for differential harnesses, based on C++, that
can integrate with the DeepSpace fuzzer and Google's FuzzTest
infrastructure, without depending on either.  This aspect of the
research will necessarily be more exploratory, initially involving
efforts to transform differential harnesses from the literature into
standalone implementations or potentially automatable AFL++ modes.

\subsection{Products}

Primary products of this proposal are therefore:

\begin{enumerate}
\item An implementation of differential-aware fuzzing algorithms as a
  new mode of operation for the AFL++ fuzzer, including novel
  seed scheduling, energy, and mutation approaches based on
  methods; these will include methods that generalized to all
  differential fuzzing problems and
  to the specialized case of versions with largely overlapping
  implementations.
  \item Modalities for automating common cases where the task is
    differential testing of two binaries that operate over a
   file or standard input, the common use of AFL and other fuzzers:
   these will include calibration tools for taking binaries and
   automatically proposing a differential comparison scheme (composed
   of those available to users) that is sufficiently sensitive to
   detect subtle bugs, but applies simple abstractions and choices of
   detection channels to reduce obvious false positives.
\item A framework/light DSL for declarative, legible, differential harness
  construction that works with the above fuzzing algorithms and cooperates
  with the DeepState unit-fuzzing framework, but does not require use
  of DeepState.
\end{enumerate}

\subsection{Evaluation}

These products will be evaluated in multiple ways.  First, algorithmic
fuzzing changes can be evaluated using field-standard methods for
evaluating fuzzing research~\cite{FuzzerHicks}.  The same differential
fuzzing problems can be executed using both baseline AFL++ and
combinations of proposed techniques (in the standard ablation study
approach) to determine the overall effectiveness of proposed methods
and of each proposed component.  Benchmarks for this aspect of
evaluation can be drawn both from existing differential harnesses
available in open source repositories and in the differential fuzzing
literature: for standard benchmarks in fuzzing literature,
differential problems can be constructed by using different software
versions with known divergences or by using source or binary mutation,
removing all mutants detected by a standard corpus (forcing
non-trivial differential problems).  For problems taken from the
differential fuzzing literature, including special-purpose
differential fuzzing for a fixed domain, numerous real-world examples
should be available, with complex problems across very different
implementations, not just differening versions of the same problem.

For these experiments studying the basic effectiveness of relational
coverage and relationally-guided mutation in fuzzing, the most
interesting benchmarks will correspond to standard fuzzing metrics:
time to first detected divergence and number of known divergences
discovered (we expect there may be tradeoffs between these, which we
aim to explore in a more systematic way than most fuzzing evaluations).

The thrust of differential specification work focusing on binary pairs
and automatic oracle generation can be evaluated by reporting false
positive and false negative rates, determining the
effectiveness of the auto-calibration modes.

Evaluating a manual harness construction DSL is less straightforward
and quantitative; we plan first to simply compare code sizes of
harnesses, based on problems from the differential fuzzing literature,
implementing the same specification with and without the DSL
constructs, which should demonstrate the ability of the approach to
concisely express standard differential testing problems.  A second
approach, tied to our expectation that differential testing may be an
important tool in AI-based code synthesis, is to pose problems in both
harness construction and harness understanding to state of the art
coding AIs, with and without the DSL.  Our hypothesis is that given a
good DSL and documentation for it, AIs will be able to construct more
effective harnesses and answer questions about harnesses better than
when differential specifications must be expressed in less concise
code, without idiomatic ``shortcuts.''

